% ---------------- Packages ----------------
%\usepackage{hyphenat}

\usepackage[utf8]{inputenc}
\usepackage{graphicx}
% \usepackage[fleqn]{amsmath}
\usepackage{amsmath}
\usepackage{listings}
\usepackage[hidelinks]{hyperref}
% \usepackage{dtsyntax}
% \usepackage[colorinlistoftodos,prependcaption,textsize=tiny]{todonotes}
\usepackage[normalem]{ulem}
\usepackage{pdflscape}
\usepackage{amssymb}
\usepackage{longtable}
%\usepackage{tikz}
%\usepackage{algorithm}
\usepackage{tikz}
\usetikzlibrary{decorations.pathmorphing, calc, arrows.meta}

\usepackage{placeins}

\usepackage[a4paper,top=2.5cm,bottom=3cm,inner=2.5cm,outer=2.5cm]{geometry}
% \usepackage[letterpaper,top=2cm,bottom=2cm,left=4cm,right=4cm,marginparwidth=1.75cm]{geometry}

\usepackage{multirow}
% For \includegraphics with valign
\usepackage[export]{adjustbox}

\usepackage{tablefootnote}

% M: Tabla en vez de Cuadro
\usepackage[spanish,es-tabla]{babel}

\usepackage{xargs} % Use more than one optional parameter in a new commands

% Algorithms boxes
\usepackage{algorithm}
\usepackage{algorithmic}

% For \qtyrange and \hertz
\usepackage{siunitx}
\usepackage{textcomp}
\DeclareSIUnit{\pixel}{px}
\DeclareSIUnit{\fps}{fps}

\usepackage{subfig}
% \usepackage{subcaption}

% Code c++
\usepackage{listings}
\usepackage{xcolor}
\usepackage{tabularx}


\renewcommand{\lstlistingname}{Código}

% Define Modelica language style for listings
\lstdefinelanguage{Modelica}{
	keywords={
		algorithm, and, annotation, block, break, class, connect, connector,
		constant, discrete, each, else, elseif, elsewhen, encapsulated, end,
		enumeration, equation, expandable, extends, external, false, final,
		flow, for, function, if, import, in, initial, inner, input, loop,
		model, not, operator, or, outer, output, package, parameter, partial,
		protected, public, record, redeclare, replaceable, return, stream,
		then, true, type, when, while, within
	},
	sensitive=true,
	comment=[l]{//},
	morecomment=[s]{/*}{*/},
	morestring=[b]",
	morestring=[b]'
}

% Configure general listing settings
\lstset{
	language=Modelica,
	basicstyle=\small\ttfamily,
	keywordstyle=\color{blue}\bfseries,
	commentstyle=\color{gray}\itshape,
	stringstyle=\color{teal},
	numbers=left,
	numberstyle=\tiny\color{gray},
	stepnumber=1,
	numbersep=8pt,
	showstringspaces=false,
	breaklines=true,
	frame=single,
	captionpos=b
}


% ---------------- Macros / Custom commands ----------------

% Title page
\newcommand*{\titleGP}{\begingroup
    \centering
    \vspace*{\baselineskip}
    \thispagestyle{empty}

    \rule{\textwidth}{1.6pt}\vspace*{-\baselineskip}\vspace*{2pt}
    \rule{\textwidth}{0.4pt}\\[\baselineskip]

    {\LARGE Simulador en tiempo real de dinámica vehicular usando modelos orientados
a objetos. } \\[0.2\baselineskip]

    \rule{\textwidth}{0.4pt}\vspace*{-\baselineskip}\vspace{3.2pt}
    \rule{\textwidth}{1.6pt}\\[\baselineskip]

    \vspace*{2\baselineskip}

    \includegraphics[width=0.3\textwidth]{img/logo-unr.png}

    \vspace*{2\baselineskip}

    \scshape
    Tesina de Grado\\[\baselineskip]
    Licenciatura en Ciencias de la Computación\par

    \vspace*{2\baselineskip}

    {\Large Autor: Juan Segundo Valero\\}
    {Legajo: V-3012/1}
    \vspace*{2\baselineskip}
    {\Large \\ Director: Ernesto Kofman} \vfill

    \vspace*{2\baselineskip}
    {\scshape \today} \\[0.3\baselineskip]

\endgroup}
